\section{Threats to Validity}
\label{sec:threats}

\paragraph{Internal Validity.}
Fuzz testing is inherently non-deterministic.
Random seeds, mutator selection choices, and system scheduling can produce variation across experimental runs.
To mitigate this threat:
(i)~For \fastgltf, we executed 5 independent trials per mutation arm with random seeds, reporting median values and interquartile ranges (IQR) across 15 trials;
(ii)~For \raylib and \drogon, we executed extensive multi-worker campaigns totaling 54.3 million and 750 million inputs to observe long-term convergence;
(iii)~We verified every discovered defect using deterministic, standalone C++ reproducers that replay the exact crashing input independently of the fuzzer runtime.

\paragraph{External Validity.}
Our evaluation examines three open-source C and C++ projects: \fastgltf, \raylib, and \drogon.
These three projects may not be representative of all C/C++ software.
To mitigate this threat, we selected mature, production-grade systems spanning three distinct architectural domains: 3D binary file parsing, interactive multimedia rendering, and high-concurrency network servers.
However, our conclusions may not directly generalize to bare-metal embedded firmware or safety-critical real-time operating systems.

\paragraph{Construct Validity.}
A potential threat is whether branch edge coverage accurately reflects software security testing effectiveness.
High coverage does not guarantee that all vulnerabilities will be triggered.
To mitigate this threat, we paired coverage feedback with AddressSanitizer and UndefinedBehaviorSanitizer instrumentation.
Furthermore, we evaluated bug-finding effectiveness by reporting candidate defects to the upstream maintainers and confirming their security impact through merged pull requests and official bug trackers.
